\section{سخت‌افزار و ساختار فایل‌ها}
این راهنما برای برد
\lr{NodeMCU 1.0 ESP-12E}
و کابل دادهٔ
\lr{USB}
نوشته شده است. پروژهٔ
\lr{firmware}
فقط هشت فایل لازم دارد:

\begin{LTR}
\begin{lstlisting}
code/arduino/mqtt/
|-- platformio.ini
|-- include/
|   |-- actuators.h
|   |-- sensors.h
|   |-- secrets.example.h
|   `-- secrets.h          # local and ignored by Git
`-- src/
    |-- actuators.cpp
    |-- main.cpp
    `-- sensors.cpp
\end{lstlisting}
\end{LTR}

\section{تنظیم اطلاعات محلی}
ابتدا وارد پوشهٔ پروژهٔ برد شوید و فایل نمونه را کپی کنید:

\begin{LTR}
\begin{lstlisting}
cd code/arduino/mqtt

# Linux / macOS
cp include/secrets.example.h include/secrets.h

# Windows PowerShell
Copy-Item include/secrets.example.h include/secrets.h
\end{lstlisting}
\end{LTR}

در فایل
\lr{include/secrets.h}
نام و رمز شبکه را وارد کنید. این فایل توسط
\lr{Git}
نادیده گرفته می‌شود. رمز را در
\lr{main.cpp}
یا مستندات ننویسید.

سپس در
\lr{src/main.cpp}
مقدار
\lr{MQTT\_HOST}
را برابر نشانی رایانه قرار دهید:

\begin{LTR}
\begin{lstlisting}[language=C++]
constexpr char MQTT_HOST[] = "(*@\lr{192.168.1.15}@*)";
\end{lstlisting}
\end{LTR}

این نشانی فقط نمونه است. نشانی
\lr{127.0.0.1}
یا
\lr{localhost}
برای برد قابل استفاده نیست.

\section{رفتار
\lr{firmware}
}
کد فعلی پنج مسئولیت دارد:

\begin{enumerate}
  \item اتصال به شبکه و نمایش نشانی برد؛
  \item اتصال به کارگزار و تلاش دوباره بدون متوقف‌کردن برنامه؛
  \item انتشار وضعیت
  \lr{online/offline}
  ؛
  \item خواندن دما و رطوبت آزمایشی از ماژول سنسورها؛
  \item ساخت و انتشار بسته‌های داده هر پنج ثانیه؛
  \item دریافت فرمان و تحویل
  \lr{action}ها
  به ماژول
  \lr{actuator}.
\end{enumerate}

مقدارهای دما و رطوبت فعلاً شبیه‌سازی می‌شوند. فایل
\lr{sensors.cpp}
تنها محل وابسته به سخت‌افزار است و
\lr{actuators.cpp}
تنها محل تنظیم خروجی‌ها و قواعد ایمنی است.
\lr{main.cpp}
بسته‌بندی و ارتباط را مدیریت می‌کند.

\section{موضوع‌ها و پیام‌ها}
شناسهٔ دستگاه
\lr{esp8266-01}
است و دو موضوع زیر استفاده می‌شود:

\begin{LTR}
\begin{lstlisting}
iot/devices/esp8266-01/status
iot/devices/esp8266-01/telemetry
iot/devices/esp8266-01/control
iot/devices/esp8266-01/response
\end{lstlisting}
\end{LTR}

نمونهٔ پیام داده:

\begin{LTR}
\begin{lstlisting}
{
  "schema_version": 1,
  "device_id": "esp8266-01",
  "message_id": 42,
  "sampled_at_ms": 15342,
  "batch_index": 1,
  "batch_count": 1,
  "readings": [
    {"sensor_id":"temperature","value":(*@\lr{24.6}@*),"unit":"celsius"},
    {"sensor_id":"humidity","value":(*@\lr{58.2}@*),"unit":"percent"}
  ]
}
\end{lstlisting}
\end{LTR}

مقدارهای کنونی ساختگی‌اند و اندازه‌گیری محیط محسوب نمی‌شوند. پیام وضعیت به‌صورت نگه‌داری‌شده منتشر می‌شود و در قطع ناگهانی، سازوکار آخرین وصیت
\lr{MQTT}
وضعیت را به
\lr{offline}
تغییر می‌دهد.

\subsection{معنای فیلدهای بسته}
\begin{itemize}
  \item
  \lr{schema\_version}
  نسخهٔ ساختار پیام است؛
  \item
  \lr{message\_id}
  شمارهٔ یک نوبت نمونه‌برداری است؛
  \item
  \lr{sampled\_at\_ms}
  زمان نمونه‌برداری از لحظهٔ روشن‌شدن برد است؛
  \item
  \lr{batch\_index}
  شمارهٔ بسته و
  \lr{batch\_count}
  تعداد کل بسته‌های همان نمونه‌اند؛
  \item
  \lr{readings}
  فهرست سنسورهاست و هر عضو شناسهٔ پایدار، مقدار عددی و واحد دارد.
\end{itemize}

\section{افزودن و افزایش تعداد سنسورها}
برای افزودن سنسور، فقط تابع
\lr{readSensors}
در فایل
\lr{src/sensors.cpp}
را تغییر دهید. نمونه:

\begin{LTR}
\begin{lstlisting}[language=C++]
namespace {
constexpr size_t CURRENT_SENSOR_COUNT = 3;
}

readings[0] = {"temperature", (*@\lr{24.6F}@*), "celsius"};
readings[1] = {"humidity", (*@\lr{58.2F}@*), "percent"};
readings[2] = {"room_2_temperature", (*@\lr{22.8F}@*), "celsius"};
\end{lstlisting}
\end{LTR}

شناسهٔ سنسور باید کوتاه، یکتا و پس از شروع ذخیرهٔ داده ثابت باشد. واحد را داخل شناسه قرار ندهید؛ واحد فیلد جداگانه دارد. همهٔ داده‌های این نسخه عددی هستند. نوع‌های متنی یا منطقی فقط زمانی به قرارداد اضافه می‌شوند که یک نیاز واقعی وجود داشته باشد.

آرایهٔ فعلی تا
\lr{128}
سنسور را نگه می‌دارد. هر بسته حداکثر شش سنسور دارد؛ بنابراین دو سنسور در یک بسته، بیست سنسور در چهار بسته و صد سنسور در هفده بسته ارسال می‌شوند. همهٔ بسته‌های یک نوبت
\lr{message\_id}
یکسان دارند و دریافت‌کننده آن‌ها را با شمارهٔ بسته مرتب می‌کند. این روش از ساخت یک پیام بسیار بزرگ و مصرف ناگهانی حافظه جلوگیری می‌کند.

برای بیش از
\lr{128}
سنسور ابتدا مقدار
\lr{MAX\_SENSOR\_COUNT}
را افزایش دهید و مصرف حافظه را پس از ساخت بررسی کنید. اندازهٔ هر بسته را بدون اندازه‌گیری افزایش ندهید.

\section{ساخت
\lr{firmware}
}
نسخهٔ محیط و کتابخانه‌ها در
\lr{platformio.ini}
ثابت شده‌اند تا ساخت روی سیستم‌های مختلف تکرارپذیر باشد. در پوشهٔ پروژه اجرا کنید:

\begin{LTR}
\begin{lstlisting}
pio run
\end{lstlisting}
\end{LTR}

خروجی موفق با عبارت زیر پایان می‌یابد:

\begin{LTR}
\begin{lstlisting}
========================= [SUCCESS] =========================
\end{lstlisting}
\end{LTR}

در آزمایش انجام‌شده برای این نسخه، ساخت موفق بود و حدود
\lr{37.5\%}
حافظهٔ
\lr{RAM}
و
\lr{27.3\%}
حافظهٔ برنامه مصرف شد.

\section{ارسال روی برد و مشاهدهٔ خروجی}
برد را با کابل داده وصل کنید و درگاه‌ها را بررسی کنید:

\begin{LTR}
\begin{lstlisting}
pio device list
\end{lstlisting}
\end{LTR}

سپس
\lr{firmware}
را ارسال و پایانهٔ سریال را باز کنید:

\begin{LTR}
\begin{lstlisting}
pio run --target upload
pio device monitor
\end{lstlisting}
\end{LTR}

خروجی مورد انتظار:

\begin{LTR}
\begin{lstlisting}
Wi-Fi connected. Device IP: (*@\lr{192.168.1.x}@*)
MQTT connected.
Telemetry batch published: {"schema_version":1,...}
\end{lstlisting}
\end{LTR}

\section{تأیید پیام روی رایانه}
در پایانهٔ دیگری همهٔ پیام‌های دستگاه را دریافت کنید:

\begin{LTR}
\begin{lstlisting}
mosquitto_sub -h (*@\lr{192.168.1.15}@*) \
  -t "iot/devices/esp8266-01/#" -v
\end{lstlisting}
\end{LTR}

نشانی نمونه را با نشانی رایانه جایگزین کنید. ابتدا باید پیام وضعیت و سپس پیام‌های دوره‌ای داده دیده شوند.

\section{آزمایش اتصال دوباره}
شبکهٔ بی‌سیم را برای مدت کوتاه قطع و دوباره وصل کنید. برنامه نباید قفل شود. پس از برگشت شبکه باید دوباره پیام
\lr{MQTT connected.}
و داده‌های دوره‌ای دیده شوند.

\section{عیب‌یابی کوتاه}
\begin{itemize}
  \item حالت اتصال
  \lr{-2}
  معمولاً به نشانی نادرست، دیوارهٔ آتش یا تنظیم
  \lr{listener}
  مربوط است.
  \item برد
  \lr{ESP8266}
  به شبکهٔ
  \lr{2.4 GHz}
  نیاز دارد.
  \item نبودن درگاه سریال معمولاً از کابل بدون مسیر داده یا درایور مبدل
  \lr{USB}
  است.
  \item اگر
  \lr{upload}
  انجام نشد، برنامه‌ای را که درگاه سریال را باز نگه داشته است ببندید و دوباره تلاش کنید.
\end{itemize}

برای این مرحله تصویر لازم نیست؛ خروجی پایانه، موضوع پیام و متن
\lr{JSON}
شاهد دقیق‌تر و قابل جست‌وجویی برای عیب‌یابی هستند.
